% Source of input/bench/cs_book/kb/assessment-integrity/assessment_and_integrity.pdf
% Written for the AutoGenBook benchmark. Original text, Apache 2.0 (repository licence).
% Build: lualatex assessment_and_integrity.tex (twice). The running header and the
% page-number footer are deliberate: the new extractor must drop them.
\documentclass[11pt,a4paper]{article}
\usepackage{fontspec}
\usepackage[margin=28mm]{geometry}
\usepackage{fancyhdr}
\usepackage{booktabs}
\usepackage{array}
\pagestyle{fancy}
\fancyhf{}
\fancyhead[L]{Assessment and Academic Integrity in the Age of Generative AI}
\fancyfoot[C]{Page \thepage}
\renewcommand{\headrulewidth}{0.4pt}
\title{Assessment and Academic Integrity in the Age of Generative AI}
\author{AutoGenBook benchmark corpus}
\date{}
\begin{document}
\maketitle
\thispagestyle{fancy}

\section{Ethical principles}

The use of generative AI in teaching should rest on five principles:
transparency, privacy, fairness, integrity and accountability. Transparency
means that teachers and students state openly when and how generative tools
were used. Privacy means that personal data of students are never entered
into a tool without a legal basis and adequate protection. Fairness requires
that no student is disadvantaged because they cannot afford a paid tool.
Integrity means that assessed work reflects the student's own learning.
Accountability means that a human remains responsible for every decision
that affects a student.

These principles apply equally to teachers. A teacher who uses a model to
draft feedback or to grade short answers remains accountable for the grade
and must be able to explain it.

\section{Data protection}

Within the European Union the processing of personal data is governed by the
General Data Protection Regulation (GDPR). Student names, identification
numbers, grades, submitted work and even writing style can be personal data.
Before student data are processed by an external AI service, the university
must have a legal basis for the processing and a data processing agreement
with the provider.

Three practical rules follow from the regulation. The first is data
minimisation: only the data that are strictly necessary for the task should
be processed. The second is anonymisation: names and identifiers should be
removed from texts before they are pasted into a tool. The third is
documentation: when a new tool will process student data at scale, a data
protection impact assessment (DPIA) should be carried out before deployment.
Enterprise licences and locally hosted models make it much easier to meet
these obligations than free consumer accounts, whose terms often allow the
provider to use conversations for further training.

\section{The EU AI Act}

The European regulation on artificial intelligence, known as the AI Act, was
adopted in 2024 and introduces obligations that depend on the risk of an AI
system. Systems with unacceptable risk are prohibited, high-risk systems must
meet strict requirements on data quality, documentation and human oversight,
and systems with limited risk are subject mainly to transparency obligations.

Education is explicitly named among the high-risk areas. AI systems used to
decide on admission to education, to evaluate learning outcomes, or to
monitor students during tests are classified as high-risk. A chat assistant
used by a teacher to draft exercises is not high-risk in itself, but the
same technology becomes high-risk when it is used to assign grades
automatically.

\section{Categories of assignments}

Instead of a single rule for the whole university, many institutions ask
teachers to classify each assignment into one of three categories and to
state the category in the assignment brief. Table~\ref{tab:categories}
summarises the scheme.

\begin{table}[h]
\centering
\begin{tabular}{>{\raggedright}p{3.2cm} >{\raggedright}p{5.6cm} >{\raggedright\arraybackslash}p{4.2cm}}
\toprule
Category & Meaning & Typical example \\
\midrule
Prohibited & Generative tools may not be used at all & Closed-book written exam \\
Permitted with declaration & Tools may be used if the student declares how & Seminar essay \\
Required & Using a tool is part of the task & Critique of an AI-generated text \\
\bottomrule
\end{tabular}
\caption{Three categories of assignments}
\label{tab:categories}
\end{table}

For the category permitted with declaration, the declaration should name the
tool, describe what it was used for, and include the key prompts. A short
statement in the syllabus explains the categories to students at the start
of the semester and prevents misunderstandings later.

\section{Citing generative AI}

When a student uses text or ideas produced by a generative tool, the use
should be acknowledged like any other source. The major citation styles now
provide templates for this purpose. In APA style, for example, the reference
names the company that developed the model as the author, the name and
version of the model, the date of use and the address of the service. The
exact prompt should be reported in the text or in an appendix, because the
same model can produce different answers to the same question.

\section{Why AI detectors are not reliable}

Software that claims to detect text written by generative AI produces both
false positives and false negatives. Independent evaluations have shown that
detectors frequently flag texts written by non-native speakers as generated,
and that light paraphrasing is enough to hide generated text from them.
A detector score is therefore not evidence of misconduct and must never be
the sole basis for an accusation. Universities that have studied the problem
recommend redesigning assessment instead of relying on detection.

\section{Redesigning assessment}

Assessment that is resistant to misuse of generative tools shifts attention
from the final product to the process of learning. Oral examinations and
short oral defences of written work allow the teacher to verify
understanding directly. Process portfolios, in which students submit drafts,
notes and reflections at several points during the semester, make it
possible to follow how the work developed. Tasks connected to local data,
personal experience or recent in-class discussion are difficult to complete
with a general-purpose model alone.

A course redesign should start from the learning outcomes: for each outcome
the teacher asks whether the current assessment still provides evidence that
the student has achieved it when generative tools are freely available. If
not, the task is changed, moved into the classroom, or complemented by an
oral component.

\section{Practical checklist}

Before the semester starts, a teacher can use the following checklist:
decide the category of every assignment; add a statement on generative AI to
the syllabus; choose tools that are covered by an institutional licence;
prepare an example of a correct declaration; and plan at least one activity
that develops the students' ability to evaluate generated content critically.
During the semester, the teacher collects feedback from students on how the
rules work in practice and adjusts them for the next run of the course.

\end{document}
